%=====================================================================
\chapter{Responsible AI and Data Ethics}\label{ch:ethics}
%=====================================================================
\locator{6}{Part VI --- Professional Practice}

\begin{outcomes}
By the end of this chapter you will be able to:
\begin{tight}
  \item \textbf{Identify} where bias enters a data science system, at each of five stages.
  \item \textbf{Compute} and \textbf{compare} the main fairness metrics, and explain why they cannot all hold at once.
  \item \textbf{Apply} the core data-protection principles to a project design.
  \item \textbf{Choose} an explainability technique appropriate to the audience and the stakes.
  \item \textbf{Conduct} a structured ethical review before deployment.
\end{tight}
\end{outcomes}

\begin{prereq}
\chref{evaluation} (metrics --- fairness metrics are confusion matrices computed per
group) and \chref{causality} (proxies and confounding).
\end{prereq}

\begin{keyterms}
algorithmic bias $\bullet$ protected attribute $\bullet$ proxy variable $\bullet$
demographic parity $\bullet$ equal opportunity $\bullet$ equalised odds $\bullet$
calibration $\bullet$ impossibility result $\bullet$ personal data $\bullet$ consent
$\bullet$ data minimisation $\bullet$ purpose limitation $\bullet$ anonymisation
$\bullet$ explainability $\bullet$ SHAP $\bullet$ LIME $\bullet$ model card
\end{keyterms}

\section{Why This Chapter Is Not Optional}

\begin{conceptbox}[The Position This Chapter Takes]
Everything in this book so far optimises a number. This chapter is about the
questions that number cannot answer: whether the system should exist, who it
affects, who can contest it, and what happens when it is wrong about a particular
person. A technically excellent system that fails these questions is a failed
system --- and increasingly, an unlawful one.
\end{conceptbox}

\section{Where Bias Enters}

\dsfig{%
\begin{tikzpicture}[font=\scriptsize,
  s/.style={rounded corners=3pt, draw=accent, fill=accent!8, text=accent!65!black,
            text width=2.35cm, align=center, font=\tiny\bfseries, minimum height=0.85cm},
  ex/.style={font=\tiny, text=gray!55!black, align=center, text width=2.6cm},
  arr/.style={-{Stealth[length=2mm]}, accent!55, line width=0.9pt}]

\node[s] (a) at (0,0)     {1. THE WORLD};
\node[s] (b) at (2.85,0)  {2. COLLECTION};
\node[s] (c) at (5.7,0)   {3. LABELS};
\node[s] (d) at (8.55,0)  {4. MODEL};
\node[s] (e) at (11.4,0)  {5. DEPLOYMENT};

\foreach \x/\y in {a/b,b/c,c/d,d/e}{\draw[arr] (\x) -- (\y);}

\node[ex] at (0,-1.35)     {Historical inequity is real. Data recording it faithfully is
                            \emph{accurate} and still unjust to learn from.};
\node[ex] at (2.85,-1.35)  {Who is missing? A VLE-only survey never reaches students who
                            stopped logging in.};
\node[ex] at (5.7,-1.35)   {``Good employee'' = past promotions, which encode past bias.
                            The label is the judgement.};
\node[ex] at (8.55,-1.35)  {Optimising overall accuracy sacrifices minority groups ---
                            they cost the loss function least.};
\node[ex] at (11.4,-1.35)  {Feedback loops: policing where you predicted crime generates
                            the data that confirms it.};

\node[anchor=north west, align=left, text width=13.6cm, font=\scriptsize, text=primary!60!black]
     at (-1.3,-3.1)
     {\textbf{Removing the protected attribute does not remove the bias.} Postcode
      predicts ethnicity; school attended predicts class; typing speed and device type
      predict disability and income. A model denied a protected attribute will
      reconstruct it from \term{proxy variables} --- and will do so \emph{silently}, so
      you cannot even measure the problem. Keep the attribute for measurement; control
      its use in the decision.};
\end{tikzpicture}}%
{Five stages at which bias enters. Only stage 4 is a modelling problem; the other four
are design, data and deployment problems, which is why fairness cannot be fixed by
choosing a different algorithm.}{bias-stages}

\section{Fairness Metrics --- and Why You Cannot Have Them All}

\begin{definitionbox}[Three Common Fairness Criteria]
\term{Demographic parity}: the flag rate is equal across groups,
$P(\hat{y}=1 \mid A{=}a)$ equal for all $a$.\\[2pt]
\term{Equal opportunity}: recall is equal across groups,
$P(\hat{y}=1 \mid y{=}1, A{=}a)$ equal for all $a$.\\[2pt]
\term{Calibration}: a score of 0.7 means the same thing in every group,
$P(y{=}1 \mid \hat{p}{=}s, A{=}a)$ equal for all $a$.
\end{definitionbox}

\begin{worked}[The Impossibility Result, in Numbers]
A dropout model is evaluated on two groups. Group A has a genuine base rate of 20\%;
group B, for reasons of prior disadvantage entirely outside the model, has 40\%.

\smallskip
\textbf{Suppose the model is well calibrated} in both groups --- a score of 0.6
means a 60\% chance of failing, whoever you are. This is usually considered the
minimum requirement for a score to be honest.

\smallskip
\textbf{Then flag rates cannot be equal.} Group B contains twice as many genuine
cases, so any calibrated model flags roughly twice as many of them. Demographic
parity is violated by construction.

\smallskip
\textbf{Force demographic parity instead} --- flag 25\% of each group. Now you must
use a lower threshold for A and a higher one for B. In group B, students at genuine
risk are left unflagged so the quota can be met; in group A, students who are fine
are flagged. The score no longer means the same thing in each group, so calibration
is broken.

\smallskip
\textbf{This is not a defect of your model.} It is a theorem: when base rates differ
between groups, calibration and equal error rates cannot both hold, except in the
degenerate case of a perfect classifier. You must \emph{choose} which fairness
criterion your context demands, state the choice explicitly, and justify it --- and
no amount of technical work will let you avoid the choice.

\smallskip
\textbf{How to choose here.} The intervention is supportive (a phone call offering
help), not punitive, and capacity is limited. Equal opportunity is the defensible
target: a student at genuine risk should have the same chance of being offered help
regardless of group. That is a value judgement, made by the institution, recorded in
writing --- not a hyperparameter.
\end{worked}

\dsfig{%
\begin{tikzpicture}
\begin{groupplot}[
  group style={group size=3 by 1, horizontal sep=0.9cm},
  width=5.35cm, height=4.0cm,
  axis lines=left, tick label style={font=\tiny}, label style={font=\tiny},
  title style={font=\scriptsize\bfseries},
  ybar={9pt}, ymin=0, ymax=0.72,
  symbolic x coords={Group A, Group B}, xtick=data,
  legend style={font=\tiny, draw=gray!40, at={(0.5,-0.22)}, anchor=north, legend columns=2},
]
\nextgroupplot[title={Calibrated model}, ylabel={rate}]
\addplot[draw=primary, fill=primary!35] coordinates {(Group A,0.20) (Group B,0.40)};
\addplot[draw=goldhl, fill=goldhl!45] coordinates {(Group A,0.22) (Group B,0.43)};
\legend{true rate, flag rate}
\node[font=\tiny, accent, align=center, anchor=north] at (axis cs:Group A,0.70)
     {flag rates differ};

\nextgroupplot[title={Forced demographic parity}]
\addplot[draw=primary, fill=primary!35] coordinates {(Group A,0.20) (Group B,0.40)};
\addplot[draw=goldhl, fill=goldhl!45] coordinates {(Group A,0.30) (Group B,0.30)};
\node[font=\tiny, accent, align=center, anchor=north] at (axis cs:Group A,0.70)
     {equal flags,\\but B under-served};

\nextgroupplot[title={Equal opportunity (recall)}]
\addplot[draw=secondary, fill=secondary!35] coordinates {(Group A,0.78) (Group B,0.77)};
\addplot[draw=accent, fill=accent!35] coordinates {(Group A,0.09) (Group B,0.11)};
\legend{recall, false-positive rate}
\node[font=\tiny, greenhl!45!black, align=center, anchor=north] at (axis cs:Group A,0.70)
     {equal chance of\\being helped};
\end{groupplot}
\end{tikzpicture}}%
{The same model under three fairness criteria. Each panel is defensible in some context
and indefensible in others. The professional obligation is to choose deliberately and say
so, not to achieve all three.}{fairness-metrics}

\section{Privacy and Data Protection}

\rowcolors{2}{white}{lightbg}
\begin{longtable}{@{}L{3.2cm}L{4.8cm}L{6.2cm}@{}}
\toprule
\rowcolor{primary!15}
\textbf{Principle} & \textbf{Means} & \textbf{In practice, for your project} \\
\midrule
\endhead
Lawful basis & You need a legal reason to process & Identify it before collection, not after the model works \\
Purpose limitation & Use it only for the stated purpose & Attendance collected for registers may not be repurposed for risk scoring without a new basis \\
Data minimisation & Collect only what you need & Every extra field is liability. Justify each one \\
Accuracy & Keep it correct and current & Stale data produces wrong decisions about real people \\
Storage limitation & Delete when no longer needed & Set retention at creation; nobody deletes later \\
Security & Protect it appropriately & Encryption, access control, audit logs (\chref{cybersecurity}) \\
Accountability & Be able to demonstrate compliance & Documentation, lineage (\chref{dataeng}), model cards \\
Rights of the individual & Access, rectification, objection, explanation & You must be able to say why \emph{this person} was scored as they were \\
\bottomrule
\end{longtable}

\begin{alertbox}[Anonymisation Is Harder Than It Looks]
Removing names is \emph{pseudonymisation}, not anonymisation. A combination of
postcode, date of birth and gender identifies most individuals uniquely; a single
student's timetable plus a few timestamps identifies them in any campus dataset.
Assume re-identification is possible unless you have specifically tested otherwise,
and treat pseudonymised data as personal data --- because legally, it usually is.
\end{alertbox}

\section{Explainability}

\rowcolors{2}{white}{defbg}
\begin{longtable}{@{}L{3.0cm}L{4.6cm}L{6.6cm}@{}}
\toprule
\rowcolor{primary!15}
\textbf{Technique} & \textbf{Gives you} & \textbf{Use when} \\
\midrule
\endhead
Use a simple model & Global, exact, free & The stakes are high and a small accuracy loss is acceptable --- often the right answer \\
Coefficients / odds ratios & Global, exact & Linear or logistic models (\chref{classification}) \\
Permutation importance & Global, model-agnostic & ``Which features matter overall?'' (\chref{features}) \\
SHAP & Per-prediction contribution of each feature, with a consistency guarantee & You must explain \emph{this} decision to \emph{this} person \\
LIME & Local linear approximation around one prediction & Quick local intuition; less stable than SHAP \\
Counterfactual & ``Had X been different, the decision would change'' & \textbf{The most actionable form} --- it tells the person what to do \\
\bottomrule
\end{longtable}

\begin{conceptbox}[Explain to the Audience, Not to Yourself]
A SHAP force plot satisfies a data scientist. A student needs: \emph{``You were
flagged mainly because you have not submitted since week 2. Submitting this week's
work would remove the flag.''} That is a counterfactual explanation, and it is the
only kind that gives the person a route to a different outcome. If your explanation
does not tell someone what they could do differently, it is documentation rather
than accountability.
\end{conceptbox}

\section{An Ethical Review You Can Actually Run}

\dsfig{%
\begin{tikzpicture}[font=\scriptsize,
  q/.style={rounded corners=3pt, draw=#1, fill=#1!7, text width=4.1cm, align=left,
            font=\tiny, inner sep=5pt, minimum height=2.5cm},
  hd/.style={rounded corners=3pt, fill=#1, text=white, font=\tiny\bfseries,
             text width=4.1cm, align=center, minimum height=0.5cm}]

\node[hd=primary] at (0,0)     {BEFORE YOU BUILD};
\node[q=primary]  at (0,-1.65) {%
  $\square$ Who is affected, including people who are not the customer?\\[3pt]
  $\square$ What is the worst outcome for the \emph{worst-affected individual}, not on
  average?\\[3pt]
  $\square$ Is there a less invasive way to achieve the same goal?\\[3pt]
  $\square$ Would we be comfortable if this were reported publicly?};

\node[hd=goldhl] at (4.75,0)     {WHILE YOU BUILD};
\node[q=goldhl]  at (4.75,-1.65) {%
  $\square$ Which groups can we measure performance for?\\[3pt]
  $\square$ Do error rates differ across them, and by how much?\\[3pt]
  $\square$ Which fairness criterion have we chosen, and who signed off?\\[3pt]
  $\square$ What proxies for protected attributes are in the feature set?};

\node[hd=greenhl] at (9.5,0)     {BEFORE YOU DEPLOY};
\node[q=greenhl]  at (9.5,-1.65) {%
  $\square$ Can we explain any single decision to the person it concerns?\\[3pt]
  $\square$ How does someone contest a decision, and who reviews it?\\[3pt]
  $\square$ What is the human's role --- decide, review, or rubber-stamp?\\[3pt]
  $\square$ What monitoring would tell us it has started harming someone?};

\node[anchor=north west, align=left, text width=13.6cm, font=\scriptsize, text=accent!75!black]
     at (-2.3,-3.5)
     {\textbf{The two questions that do most of the work} are ``what is the worst outcome
      for the worst-affected individual?'' --- because averages hide exactly the harm you
      need to see --- and ``how does someone contest this?'' --- because a system with no
      appeal route is not a decision-support tool, it is an unaccountable authority.};
\end{tikzpicture}}%
{A three-stage ethical review. It takes about an hour, requires no specialist training,
and catches most of what goes wrong.}{ethics-review}

\begin{definitionbox}[Model Card]
A short standard document accompanying every deployed model: its intended use,
out-of-scope uses, training data and its limitations, performance overall
\emph{and by subgroup}, ethical considerations, and contact details for challenge.
It is the accountability artefact that makes the rest of this chapter auditable.
\end{definitionbox}

%---------------------------------------------------------------------
\begin{fourlenses}{Responsible Practice}
\lensLA{The highest-stakes lens in this book, because the subjects are identifiable
young people with limited power in the relationship. \emph{Specific risks:}
self-fulfilling prophecy --- a student told they are ``high risk'' may disengage
further; proxy discrimination via postcode, school and device type; and mission
creep from support into assessment. \emph{Requirements:} equal opportunity as the
fairness criterion, counterfactual explanations given to the student, a human
adviser who can and does overrule the model, and a stated route to contest. Never
show a raw risk score to a student.}

\lensPM{Predicting which teams or individuals will underperform edges quickly into
workplace surveillance. \emph{Rule of thumb:} model the \emph{work}, not the
\emph{worker}. Forecasting sprint completion is legitimate; ranking developers by
predicted productivity is a different activity with employment-law consequences and
predictable effects on behaviour --- people optimise whatever is measured, usually
at the expense of what is not.}

\lensIOT{Sensing is easy to deploy and hard to consent to. Occupancy and movement
data is personal data the moment it can be linked to an individual, which is more
often than designers assume. \emph{Practices:} aggregate at the edge so that raw
identifiable data never leaves the device (\chref{cloudedge}), set retention short
and enforce it, and post clear notice where sensing occurs. ``We only store counts''
is only true if you engineered it to be.}

\lensSEC{Security work involves monitoring colleagues, which is legitimate and
requires limits. \emph{Practices:} proportionality --- monitor what the threat
justifies, not what is technically possible; access to raw logs restricted and
itself audited; and a documented process for when monitoring surfaces something
non-security-related. Note also the dual-use problem: the techniques in
\chref{cybersecurity} defend systems and also enable surveillance, and the
difference lies entirely in governance rather than in the code.}
\end{fourlenses}

%---------------------------------------------------------------------
\begin{lab}[Measuring Fairness]
\begin{lstlisting}[language=Python]
import numpy as np, pandas as pd
from sklearn.metrics import confusion_matrix

def fairness_report(y_true, y_pred, y_score, group):
    """Fairness metrics are just per-group confusion matrices (Chapter 14)."""
    rows = []
    for g in pd.unique(group):
        m = group == g
        tn, fp, fn, tp = confusion_matrix(y_true[m], y_pred[m]).ravel()
        rows.append({
            "group":      g,
            "n":          int(m.sum()),
            "base_rate":  round(y_true[m].mean(), 3),
            "flag_rate":  round(y_pred[m].mean(), 3),   # demographic parity
            "recall":     round(tp / (tp + fn), 3),     # equal opportunity
            "precision":  round(tp / (tp + fp), 3),
            "fpr":        round(fp / (fp + tn), 3),     # equalised odds (with recall)
        })
    df = pd.DataFrame(rows)

    print(df.to_string(index=False))
    print(f"\nmax recall gap    : {df.recall.max() - df.recall.min():.3f}")
    print(f"max flag-rate gap : {df.flag_rate.max() - df.flag_rate.min():.3f}")
    print("NOTE: if base rates differ, these two gaps CANNOT both be zero.")
    return df

# --- calibration check: does 0.7 mean 0.7 for everyone? ---------------
def calibration_by_group(y_true, y_score, group, bins=5):
    q = pd.qcut(y_score, bins, labels=False, duplicates="drop")
    tab = (pd.DataFrame({"bin": q, "y": y_true, "g": group})
             .groupby(["g", "bin"])["y"].mean().unstack())
    print("\nobserved rate per score bin, per group:")
    print(tab.round(3))     # rows should be similar; large differences = miscalibration

# --- proxy detection: can group be predicted from the features? -------
from sklearn.ensemble import RandomForestClassifier
from sklearn.model_selection import cross_val_score
proxy_auc = cross_val_score(RandomForestClassifier(n_estimators=200),
                            X, group, cv=5, scoring="roc_auc").mean()
print(f"\ngroup predictable from features: AUC {proxy_auc:.3f}")
# AUC well above 0.5 means proxies exist -- dropping the protected attribute
# achieved nothing except making the bias unmeasurable.
\end{lstlisting}
\end{lab}

%---------------------------------------------------------------------
\begin{checkpoint}
\begin{tightnum}
  \item You remove \texttt{gender} from the feature set. Has the model become fair?
        What test would you run to find out?
  \item Two groups have base rates of 15\% and 35\%. Can a calibrated model also
        achieve demographic parity? Explain.
  \item A student asks why they were flagged. Which explainability technique gives
        them something they can act on, and what does the explanation look like?
\end{tightnum}
\end{checkpoint}

%---------------------------------------------------------------------
\begin{chaptersummary}
\begin{tight}
  \item Bias enters at five stages; only one is a modelling problem, so no algorithm
        choice fixes it.
  \item Removing a protected attribute does not remove bias --- proxies reconstruct
        it and make it unmeasurable. Retain it for measurement, control its use.
  \item Demographic parity, equal opportunity and calibration cannot all hold when
        base rates differ. Choose deliberately, document the choice, and have it
        signed off.
  \item Data protection principles --- lawful basis, purpose limitation,
        minimisation, retention, security, accountability, individual rights ---
        constrain project design from the start.
  \item Pseudonymisation is not anonymisation; assume re-identification is possible.
  \item Explain to the audience. Counterfactual explanations are the only kind that
        give a person a route to a different outcome.
  \item Run the three-stage review, and publish a model card.
\end{tight}
\end{chaptersummary}

\begin{reviewq}
  \item \textbf{(MCQ)} Equal opportunity requires equality across groups of:
        (a)~flag rate (b)~recall (c)~precision (d)~accuracy.
  \item \textbf{(MCQ)} Postcode used as a feature is problematic mainly because it is:
        (a)~high-cardinality (b)~a proxy for protected attributes (c)~missing often
        (d)~nominal.
  \item \textbf{(Short)} State the impossibility result in your own words and explain
        what it obliges a practitioner to do.
  \item \textbf{(Short)} Distinguish pseudonymisation from anonymisation, with an
        example of re-identification.
  \item \textbf{(Short)} Why is a counterfactual explanation more useful to an
        affected individual than a feature-importance chart?
  \item \textbf{(Applied)} Write a model card for the dropout-prediction model used
        throughout this book: intended use, out-of-scope use, data, subgroup
        performance you would report, limitations, and the contest route.
  \item \textbf{(Applied)} A university proposes webcam-based exam proctoring with
        automated cheating detection. Run the three-stage ethical review of
        \dsref{ethics-review} and give a recommendation with conditions.
\end{reviewq}
